%%% FIELD dense-statement
Suppose that there is a constant \(\varepsilon\) such that
\[
c_{\rm link}<\varepsilon<C_{\rm link},
\qquad
2\varepsilon\sum_{j=2}^{\infty}j^{-a_L}<1.
\tag{1}
\]
Then the handle \(\mathsf H_{\rm law}\) has the following explicit legal realization.  For \(a\in[a_L,a_U]\), let
\[
q_a(x,w)=K^{-2}\left\{1+\varepsilon\sum_{j=2}^Kj^{-a}u_j(x)v_j(w)\right\},
\quad
J_K=\lfloor K/2\rfloor+1,
\quad
B_K=\{J_K,\ldots,K\},
\quad
M_K=|B_K|.
\tag{2}
\]
For \(s=(s_j)_{j\in B_K}\in\{-1,1\}^{M_K}\) and \(\delta>0\), put
\[
f_s(x)=\frac{\delta}{\sqrt{M_K}}\sum_{j\in B_K}s_ju_j(x),
\tag{3}
\]
and define \(P_0\) and \(P_s\) by the common \((X,W)\) mass function \(q_a\) and the conditional success probabilities
\[
P_0(Y=1\mid X=x,W=w)=\eta,
\qquad
P_s(Y=1\mid X=x,W=w)=\eta+f_s(x).
\tag{4}
\]
If \(\delta\sqrt{2M_K}\le \eta/2\), there is a constant \(c>0\), independent of \(a,K,n\), such that \(P_0,P_s\in\mathcal P_K(a)\), every one of their \((X,W,Y)\) cells has mass at least \(cK^{-2}\), and their link inequalities retain slack bounded away from zero uniformly in \(a\) and \(K\).  The null bridge \(h_0\equiv\eta\) is a regular lower-box vertex of \(\mathcal H_\beta(R,\eta)\), \(d(P_0)=0\), and every \(s\) satisfies
\[
d(P_s)\ge \delta-\varepsilon R J_K^{-(a+\beta)}.
\tag{5}
\]
Moreover, for \(\overline P_n=2^{-M_K}\sum_sP_s^{\otimes n}\),
\[
\chi^2\!\left(\overline P_n,P_0^{\otimes n}\right)
\le
\exp\!\left\{\frac{n^2\delta^4}{2M_K\eta^2(1-\eta)^2}\right\}-1.
\tag{6}
\]
Consequently, there are constants \(c_0,c_1>0\), independent of \(a,K,n\), such that, whenever \(K=K_n=o(n^{2/3})\) and
\[
K^{-(a_L+\beta)}\le c_0\frac{K^{1/4}}{\sqrt n},
\tag{7}
\]
then, for all sufficiently large \(n\), uniformly over \(a\in[a_L,a_U]\),
\[
\rho^\star_{n,K}(a,\epsilon_{\rm tot})
\ge c_1\frac{K^{1/4}}{\sqrt n}.
\tag{8}
\]
Thus the open question has a partial affirmative answer on the explicit compatibility subregime (1).  In particular, the tail restriction is imposed at the worst-case endpoint \(a_L\), and the all-sign positive construction covers every polynomial growth law \(K_n\asymp n^\kappa\) with \(\kappa<2/3\).  No claim is made here outside (1), or that \(2/3\) is an impossibility threshold for other constructions.
%%% FIELD dense-proof
Because \(a_L>1\), the series in (1) is finite.  Fix an \(\varepsilon\) satisfying (1), and then fix \(a\in[a_L,a_U]\) and \(K\ge2\).  The definitions of the DCT bases give
\[
K^{-1}\sum_{x=1}^Ku_j(x)u_\ell(x)=\mathbf 1\{j=\ell\},
\qquad
K^{-1}\sum_{w=1}^Kv_j(w)v_\ell(w)=\mathbf 1\{j=\ell\},
\tag{9}
\]
with \(u_1=v_1=1\), and
\[
|u_j(x)|\le\sqrt2,
\qquad
|v_j(w)|\le\sqrt2
\quad (j\ge2).
\tag{10}
\]
In particular, every nonconstant basis vector has uniform-coordinate mean zero.  Summing (2) in either coordinate and applying (9) therefore gives
\[
\sum_wq_a(x,w)=K^{-1},
\qquad
\sum_xq_a(x,w)=K^{-1},
\qquad
\sum_{x,w}q_a(x,w)=1.
\tag{11}
\]
By (10), \(a\ge a_L\), and (1),
\[
q_a(x,w)
\ge K^{-2}\left\{1-2\varepsilon\sum_{j=2}^Kj^{-a}\right\}
\ge c_qK^{-2},
\qquad
c_q:=1-2\varepsilon\sum_{j=2}^{\infty}j^{-a_L}>0.
\tag{12}
\]
Thus \(q_a\) is a strictly positive probability mass function with uniform marginals.  Since \(0<\eta<1/2\),
\[
\frac{\eta}{K}\le\frac1K\le\frac1{\eta K},
\tag{13}
\]
so all four marginal-balance requirements in \(\mathrm{def:model\mbox{-}class}\) hold.

Dividing (2) by the marginal \(p_X(x)=K^{-1}\), using (9), and recalling the definition of the conditional-expectation operator yields
\[
Av_1=u_1,
\qquad
Av_j=\varepsilon j^{-a}u_j
\quad (2\le j\le K).
\tag{14}
\]
Write \(g=\sum_{j=1}^Kg_jv_j\).  Under the uniform \(W\)-marginal, the restriction \(E_P\{g(W)\}=0\) is equivalent to \(g_1=0\).  Hence (9), (14), and the definition \(Lv_j=jv_j\) give
\[
\|Ag\|_{L^2_X}^2
=\varepsilon^2\sum_{j=2}^Kj^{-2a}g_j^2
=\varepsilon^2\|L^{-a}g\|_{\ell^2_{K,W}}^2.
\tag{15}
\]
It follows from (1) that the lower and upper link inequalities in \(\mathrm{def:model\mbox{-}class}\) hold with the uniform positive slack
\[
s_{\rm link}:=\min\{\varepsilon-c_{\rm link},C_{\rm link}-\varepsilon\}>0.
\tag{16}
\]

The only nonzero DCT coefficient of \(h_0\equiv\eta\) is \(\theta_1=\eta\).  The declared ranges \(0<\eta<1/2\) and \(R>1/2\) imply
\[
\sum_{j=1}^Kj^{2\beta}\theta_j^2=\eta^2<R^2,
\qquad
\eta=h_0(w)<1-\eta
\quad\text{for every }w.
\tag{17}
\]
Thus \(h_0\in\mathcal H_\beta(R,\eta)\), its source and upper-box constraints are slack, and all of its lower-box constraints are active.  To verify regularity, let \(V\) be the \(K\times K\) matrix with entries \(V_{wj}=v_j(w)\).  Equation (9) says \(K^{-1}V^\top V=I_K\), so the rows of \(V\) are linearly independent.  These rows, up to sign, are precisely the gradients in \(\theta\) of the active constraints \(\eta-\sum_j\theta_jv_j(w)\le0\).  Hence the active gradients satisfy the linear-independence constraint qualification, and \(h_0\) is the specified regular lower-box vertex.  Since conditional expectation preserves constants, (14) also gives
\[
Ah_0=\eta=b_0,
\qquad
d(P_0)=0
\tag{18}
\]
by \(\mathrm{def:bridge\mbox{-}class}\) and the definition of the observable defect.

From (3) and (10),
\[
\|f_s\|_\infty
\le \frac{\delta}{\sqrt{M_K}}\sum_{j\in B_K}|u_j(x)|
\le\delta\sqrt{2M_K}.
\tag{19}
\]
Under the stated amplitude restriction, (19) and (4) imply
\[
\frac\eta2\le P_s(Y=1\mid X=x,W=w)\le\frac{3\eta}{2},
\qquad
P_s(Y=0\mid X=x,W=w)\ge1-\frac{3\eta}{2}>\frac14.
\tag{20}
\]
The two conditional outcome probabilities under \(P_s\), and also under \(P_0\), are therefore at least \(\eta/2\).  Combining (12) and (20), every three-way cell under every displayed law is bounded below by
\[
cK^{-2},
\qquad
c:=c_q\eta/2>0.
\tag{21}
\]
Changing the outcome kernel from (4) does not alter \(q_a\), its marginals, or \(A\).  Equations (11)--(16) consequently verify directly, rather than by a preservation assertion about an unspecified subclass, that
\[
P_0,P_s\in\mathcal P_K(a).
\tag{22}
\]

Let \(\Pi_{B_K}\) be the \(L^2(P_X)\)-orthogonal projection onto \(\operatorname{span}\{u_j:j\in B_K\}\).  Since
\[
K/2\le J_K\le K,
\qquad
K/2\le M_K\le K,
\tag{23}
\]
both \(J_K\) and \(M_K\) are comparable to \(K\).  For any \(h_\theta\in\mathcal H_\beta(R,\eta)\), (14), the source constraint in \(\mathrm{def:bridge\mbox{-}class}\), and \(j\ge J_K\) on the block give
\[
\begin{aligned}
\|\Pi_{B_K}Ah_\theta\|_{L^2_X}^2
&=\varepsilon^2\sum_{j\in B_K}j^{-2a}\theta_j^2 \\
&=\varepsilon^2\sum_{j\in B_K}j^{-2(a+\beta)}j^{2\beta}\theta_j^2 \\
&\le\varepsilon^2J_K^{-2(a+\beta)}
       \sum_{j\in B_K}j^{2\beta}\theta_j^2 \\
&\le\varepsilon^2R^2J_K^{-2(a+\beta)}.
\end{aligned}
\tag{24}
\]
Under \(P_s\), (4) gives \(b_s=\eta+f_s\), so \(\Pi_{B_K}b_s=f_s\).  DCT orthonormality in (9) gives \(\|f_s\|_{L^2_X}=\delta\).  Orthogonal projection is a contraction, and the reverse triangle inequality and (24) therefore imply
\[
\begin{aligned}
d(P_s)
&=\inf_{h_\theta\in\mathcal H_\beta(R,\eta)}
  \|b_s-Ah_\theta\|_{L^2_X} \\
&\ge\inf_{h_\theta\in\mathcal H_\beta(R,\eta)}
  \|f_s-\Pi_{B_K}Ah_\theta\|_{L^2_X} \\
&\ge\delta-\varepsilon R J_K^{-(a+\beta)}.
\end{aligned}
\tag{25}
\]
This proves (5) after profiling over the complete moving source-and-box image; no inactive box constraint has been discarded from the infimum.

Write \(E_0\) for expectation under \(P_0\), and use
\(\chi^2(Q,P)=E_P\{(dQ/dP-1)^2\}\) for \(Q\ll P\).  For one observation, the likelihood ratio of \(P_s\) with respect to \(P_0\) is
\[
\ell_s(O)
=1+\frac{f_s(X)(Y-\eta)}{\eta(1-\eta)}.
\tag{26}
\]
Conditioning on \(X\), using the Bernoulli variance \(\eta(1-\eta)\) under \(P_0\), and then applying uniformity of \(P_X\) and (9), we obtain for any sign vectors \(s,t\)
\[
\begin{aligned}
E_0(\ell_s\ell_t)
&=1+\frac{\langle f_s,f_t\rangle_{L^2_X}}{\eta(1-\eta)} \\
&=1+\frac{\delta^2}{M_K\eta(1-\eta)}
      \sum_{j\in B_K}s_jt_j.
\end{aligned}
\tag{27}
\]
The amplitude restriction entails \(\delta^2\le\eta^2/(8M_K)<\eta(1-\eta)\), so the last expression in (27) is positive.  By the i.i.d. sampling assumption and the exact second-moment identity for this finite likelihood-ratio mixture,
\[
\begin{aligned}
1+\chi^2(\overline P_n,P_0^{\otimes n})
&=2^{-2M_K}\sum_{s,t}\{E_0(\ell_s\ell_t)\}^n \\
&\le2^{-2M_K}\sum_{s,t}
\exp\!\left\{
 \frac{n\delta^2}{M_K\eta(1-\eta)}
 \sum_{j\in B_K}s_jt_j
\right\} \\
&=\cosh\!\left(
 \frac{n\delta^2}{M_K\eta(1-\eta)}
\right)^{M_K}.
\end{aligned}
\tag{28}
\]
Here the inequality is \((1+x)^n\le e^{nx}\) for \(x>-1\), and the final equality follows by averaging the independent Rademacher products \(s_jt_j\).  Since \((\log\cosh x)'=\tanh x\le x\) for \(x\ge0\), integration from zero and evenness give \(\log\cosh x\le x^2/2\) for every real \(x\).  Applying this bound to (28) proves (6).

It remains to select the amplitude uniformly.  Choose \(c_*>0\), depending only on \(\eta\) and \(\epsilon_{\rm tot}\), so small that
\[
\frac12\left[
\exp\!\left\{\frac{c_*^4}{2\eta^2(1-\eta)^2}\right\}-1
\right]^{1/2}<1-\epsilon_{\rm tot},
\tag{29}
\]
and set
\[
\delta=c_*\frac{M_K^{1/4}}{\sqrt n}.
\tag{30}
\]
By (23), if \(K=o(n^{2/3})\), then
\[
\delta\sqrt{2M_K}
\le\sqrt2c_*\frac{K^{3/4}}{\sqrt n}=o(1),
\tag{31}
\]
so the amplitude condition holds for all sufficiently large \(n\).  This also shows directly that every polynomial law \(K_n\asymp n^\kappa\) with \(\kappa<2/3\) is covered by this positivity argument.

For \(a\in[a_L,a_U]\), (23) and \(K\ge2\) imply
\[
J_K^{-(a+\beta)}
\le2^{a_U+\beta}K^{-(a_L+\beta)}.
\tag{32}
\]
Choose \(c_0>0\) sufficiently small that
\[
\varepsilon R2^{a_U+\beta}c_0
\le c_*2^{-5/4}.
\tag{33}
\]
Under (7), equations (23), (30), (32), and (33) yield
\[
\varepsilon RJ_K^{-(a+\beta)}\le\frac\delta2,
\qquad
d(P_s)\ge\frac\delta2
\ge c_1\frac{K^{1/4}}{\sqrt n},
\qquad
c_1:=c_*2^{-5/4}.
\tag{34}
\]
This proves explicitly that the restriction at \(a_L\) controls the image tail for every \(a\ge a_L\).

Substitution of (30) into (6) makes its exponent equal to \(c_*^4/\{2\eta^2(1-\eta)^2\}\).  For any \(Q\ll P\), Cauchy--Schwarz applied to \(dQ/dP-1\) gives
\[
\|Q-P\|_{\rm TV}
=\frac12E_P\left|\frac{dQ}{dP}-1\right|
\le\frac12\sqrt{\chi^2(Q,P)}.
\tag{35}
\]
Thus (6), (29), and (35) show
\[
\|\overline P_n-P_0^{\otimes n}\|_{\rm TV}<1-\epsilon_{\rm tot}.
\tag{36}
\]
For every measurable test \(\phi=\phi(O_1,\ldots,O_n)\in[0,1]\), the variational definition of total variation, followed by the fact that a supremum dominates a finite average, gives
\[
\begin{aligned}
&\sup_{P\in\mathcal P_K(a):d(P)=0}E_{P^{\otimes n}}\phi
+\sup_{P\in\mathcal P_K(a):d(P)\ge c_1K^{1/4}/\sqrt n}
 E_{P^{\otimes n}}(1-\phi) \\
&\quad\ge
E_{P_0^{\otimes n}}\phi
+2^{-M_K}\sum_sE_{P_s^{\otimes n}}(1-\phi) \\
&\quad\ge1-\|\overline P_n-P_0^{\otimes n}\|_{\rm TV}
>\epsilon_{\rm tot}.
\end{aligned}
\tag{37}
\]
The null and alternatives in (37) are legal by (18), (22), and (34).  Hence no test satisfies the total-error requirement at this separation.  By \(\mathrm{def:minimax\mbox{-}radius}\), and monotonicity of its alternative set in the separation, (8) follows.  All constants used in (8) depend only on the fixed quantities in (1), \(a_L,a_U,\beta,R,\eta\), and \(\epsilon_{\rm tot}\), and not on \(a,K,n\).
